%------------------------------------
%           config
%------------------------------------
% Hier steht, was auf dem Deckblatt steht und welche Teile
% der Arbeit gedruckt werden. Mögliche Werte für Schalter
% und Abschlüsse stehen in characteristics.tex.
% PDF-Titel und Autor für Texifier stehen in texifier.tex.
%
% Schalter: \show ein, \hide aus.
%------------------------------------


%------------------------------------
%           Was gedruckt wird
%------------------------------------
% JKU-Deckblatt.
\def\showJKUCover			{\show}

% Englisches Abstract.
\def\showAbstract			{\show}

% Deutsche Kurzfassung. Die Arbeit selbst bleibt Englisch.
\def\showAbstractBoth		{\show}

% Inhaltsverzeichnis.
\def\showListOfContent		{\show}

% Abbildungsverzeichnis, im Literaturteil.
\def\showListOfFigures		{\show}

% Tabellenverzeichnis, im Literaturteil.
\def\showListOfTables		{\show}

% Listingverzeichnis, im Literaturteil.
\def\showListOfSourceCode	{\show}

% Literaturverzeichnis aus references/references.bib.
\def\showReferences			{\show}

% Eidesstattliche Erklärung, englisch und deutsch.
\def\showDeclaration		{\show}

% "Submitted by" und der Name auf dem Deckblatt.
\def\showSubmitter			{\show}

% Betreuer auf dem Deckblatt. Der Name steht in \supervisor.
% Bei einer Dissertation lautet die Zeile "Supervisor and First Examiner".
\def\showSupervisor			{\show}

% "Submitted at" auf dem Deckblatt. Der Text steht in \submittedAt.
\def\showSubmittedAt		{\show}

% Datum auf dem Deckblatt. Der Text steht in \date.
\def\showDate				{\show}


%------------------------------------
%           Satz
%------------------------------------
% \book    Doppelseitig. Kapitel beginnen auf einer rechten Seite.
%          Dazwischen liegt bei Bedarf eine leere linke Seite.
% \normal  Einseitig. Keine leeren Rückseiten.
%          Ein Kapitel beginnt weiterhin auf einer neuen Seite.
% Die Ränder sind in beiden Fällen gleich.
\def\printLayout			{\book}


%------------------------------------
%           Art der Arbeit
%------------------------------------
% \DisThesis, \DipThesis, \MasThesis oder \BacThesis.
% Steuert die Bezeichnung auf dem Deckblatt und in der Erklärung.
\def\type					{\MasThesis}


%------------------------------------
%           Personen
%------------------------------------
% Reihenfolge auf dem Deckblatt: Nachname Vorname.
\def\name					{Author Name}

% Erstbetreuer. \\ bricht die Zeile auf dem Deckblatt um.
% Leer lassen, wenn \showSupervisor auf \hide steht.
\def\supervisor				{Supervisor Name}

% Zweitbetreuer, nur bei einer Dissertation. Leer = nicht anzeigen.
\def\secondexaminer			{}

% Mitbetreuung. Leer = nicht anzeigen.
\def\assist					{}

% Firmenbetreuer. Beide leer lassen, wenn es keinen gibt.
% \company ist die Firma, \supervisorComp der Name.
\def\company				{}
\def\supervisorComp			{}


%------------------------------------
%           Abschluss und Studium
%------------------------------------
% \degree: \DegreeMSc, \DegreeBSc, \DegreeDI, \DegreeDrTW, \DegreeMag, ...
% \study:  \CompSc, \AI, \Inf, \TechnischeW, ...
% Die ausgeschriebenen Texte stehen in characteristics.tex.
\def\degree					{Degree}
\def\study					{Study Program}


%------------------------------------
%           Deckblatt
%------------------------------------
% Titel auf dem Deckblatt. \\ beginnt eine neue Zeile.
% Der einzeilige PDF-Titel steht in texifier.tex.
\def\CoverTitle				{Thesis Title}

% Datum auf dem Deckblatt, freier Text.
\def\date					{Month Year}

% Institut unter "Submitted at". \\ beginnt eine neue Zeile.
\def\submittedAt			{Institute of\newline System Software}
